\section*{CARTAN SUBALGEBRAS}

\begin{enumerate}
\def\labelenumi{\arabic{enumi})}
\item
  Let E be the set of commutative subalgebras of g that are reductive in g; this is also the set of commutative subalgebras of g all of whose elements are semi-simple. The Cartan subalgebras of g are the maximal elements of E.
\item
  Let \(x\) be a regular element of \(\mathfrak{g}\) . Then \(x\) is semi-simple. There exists a unique Cartan subalgebra of \(\mathfrak{g}\) containing \(x\) ; it is the commutant of \(x\) in \(\mathfrak{g}\) .
\item
  Let \(x\) be a semi-simple element of \(\mathfrak{g}\) . Then \(x\) belongs to a Cartan subalgebra of \(\mathfrak{g}\) . The element \(x\) is regular if and only if the dimension of the commutant of \(x\) is equal to the rank of \(\mathfrak{g}\) .
\item
  Let \(\mathfrak{h}\) be a Cartan subalgebra of \(\mathfrak{g}\) . Then \(\mathfrak{h}\) is said to be splitting if \(\operatorname{ad}_{\mathfrak{g}}x\) is triangularizable for all \(x \in \mathfrak{h}\) . Moreover, \(\mathfrak{g}\) is said to be splittable if \(\mathfrak{g}\) has a splitting Cartan subalgebra (this is the case if \(k\) is algebraically closed). A split semi-simple Lie algebra is a pair \((\mathfrak{g}, \mathfrak{h})\) where \(\mathfrak{g}\) is a semi-simple Lie algebra and \(\mathfrak{h}\) is a splitting Cartan subalgebra of \(\mathfrak{g}\) .
\end{enumerate}

In the remainder of this summary, \((\mathfrak{g},\mathfrak{h})\) denotes a split semi-simple Lie algebra.

